\chapter{The Arbitrage Zoo}
\label{ch:arbitrage_zoo}
\chaptermark{The arbitrage zoo}

This chapter tours the zoo of textbook arbitrages: the strategies
a quant looks for first when evaluating a new market. Each
follows the template \emph{observable}, \emph{trade},
\emph{P\&L}, \emph{why it usually doesn't exist}, \emph{when it
does}. The goal is to fix that template so firmly that, given any
market data dump, the right diagnostic question fires
automatically: when you see this shape in a market, set up this
trade.

Seven arbitrages are presented in the same fixed format, two to
three pages each. Once the template is internalised, encountering
an unfamiliar instrument (a new commodity future, a token on
an exotic crypto venue, a thinly-traded ETF) becomes a matter
of ticking through the seven types to see which, if any, is open.
An eighth section closes the chapter by stating that statistical
arbitrage (\cref{ch:mean_reversion_ou},
\cref{ch:pairs_cointegration}) is \emph{not} a member of the zoo:
it shares the name but not the no-risk guarantee, and conflating
the two has wiped out more firms than the other six combined.

\begin{backgroundbox}[Prerequisites]
From earlier:
\begin{itemize}
  \item \Cref{ch:lob}: the limit order book, the bid--ask
  spread, and \emph{walking the book}. Several arbitrages here
  are nothing more than walking two books simultaneously, one as
  buyer and one as seller.
  \item \Cref{ch:basket_etf_arb}: the ETF creation / redemption
  mechanism and authorised participants. The ETF c/r entry
  refers back rather than re-deriving.
  \item Forward: \cref{ch:options_payoffs} for European call/put
  payoffs. The put-call parity entry uses only the payoff
  diagrams; readers unfamiliar with $\Call$, $\Put$, $\Spx$,
  $\Kstrike$ should skim that section's opener first.
  \item Discounting: a riskless cashflow $X$ at time $T$ is
  worth $X e^{-\rrf T}$ today, with $\rrf$ the risk-free rate.
  No stochastic calculus is used; every payoff is deterministic
  by construction.
\end{itemize}
No prior finance vocabulary is assumed beyond \cref{ch:lob}
through \cref{ch:momentum}.
\end{backgroundbox}

\section{The arbitrage template}
\label{sec:ch13_template}

We fix the template first. The seven case studies that follow
are instances of the same five-step pattern.

\begin{definition}[Arbitrage]
\index{arbitrage|textbf}
\label{def:arbitrage}
An \emph{arbitrage} is a trading strategy that, executed at
market prices, has zero or negative initial cost, non-negative
payoff in every state of the world, and strictly positive payoff
in at least one state. Equivalently: a portfolio that costs
nothing (or pays you to hold it) and cannot lose money. The
Type~I / Type~II split in the technical literature does not
matter operationally here.
\end{definition}

``In every state of the world'' does all the work. A strategy
that wins on \emph{average} but can lose is not an arbitrage; it
is at best a positive-EV bet, requiring the sizing and risk
apparatus of \cref{ch:kelly_risk}. A true arbitrage is a free
lunch, and the central fact of mature markets is that free
lunches are rare and short-lived. This chapter is about
recognising the few that exist.

Every arbitrage in the zoo follows the same five-step pattern.

\begin{keyfactbox}[The arbitrage template]
\label{kf:arb_template}
For any candidate arbitrage, answer the following five questions
in order. If any answer is unsatisfactory, the arbitrage is not
open and you move on.
\begin{enumerate}
  \item \textbf{Observable}: What do you see in market data that
  tells you the arbitrage is open? This is the trigger
  condition: a numerical inequality between observable prices.
  \item \textbf{Trade}: What is the explicit set of orders that
  captures the discrepancy? Specify each leg, on which venue,
  in which direction, in what size.
  \item \textbf{P\&L}: What is the \emph{deterministic} profit
  in the ideal case (frictionless, fully filled, instant
  execution)? This is a number, not a distribution.
  \item \textbf{Why it usually doesn't exist}: Which forces in a
  normal, healthy market keep the inequality from opening?
  Knowing this is what tells you when the arb is real and when
  it is a data artefact or a stale quote.
  \item \textbf{When it does}: What specific conditions
  temporarily disable those forces and let the arbitrage open?
  This is the catalogue entry: the situations in which to
  watch.
\end{enumerate}
\end{keyfactbox}

\begin{intuitionbox}
The template does two jobs. Steps~1--3 are the trade: see the
gap, do the orders, count the money. Steps~4--5 are the
\emph{theory}: an arbitrage is interesting only insofar as it is
rare, because anything common is already arbed out by someone
faster. Winning relative-value desks know the closures and
openings of each arbitrage so intimately that they are
positioned for the open before it is visible to everyone else.
\end{intuitionbox}

In steady state every arbitrage here is closed, because if it
were not, someone would close it. Capturing an open arb is
mechanical once you know it is open; the interesting question
is what would have to break for the arb to open, and how you
would know. Each entry answers both.

\section{Triangular FX arbitrage}
\label{sec:ch13_triangular}

Currency pairs are quoted bilaterally (EUR/USD, USD/JPY,
EUR/JPY), and the three prices are linked by a consistency
requirement: in a frictionless market, walking the triangle EUR
$\to$ USD $\to$ JPY $\to$ EUR returns exactly one EUR. A round
trip above one is an arbitrage; below one, the reverse direction
is.

\textbf{Observable.} Three FX rates whose product around a
triangle differs from one. Concretely, denote the price of one
unit of currency~$X$ in units of currency~$Y$ by $E_{X/Y}$. The
arbitrage indicator is
\begin{equation}
  \Delta \;\defeq\; E_{\text{EUR}/\text{USD}} \cdot
  E_{\text{USD}/\text{JPY}} \cdot E_{\text{JPY}/\text{EUR}}
  \;-\; 1.
  \label{eq:triangular_indicator}
\end{equation}
If $\Delta > 0$, the EUR $\to$ USD $\to$ JPY $\to$ EUR direction
is profitable; if $\Delta < 0$, the reverse direction is. If
$\Delta = 0$ the triangle is consistent and no arbitrage exists.

\textbf{Trade.} Starting with one EUR:
\begin{enumerate}
  \item Sell EUR for USD at rate $E_{\text{EUR}/\text{USD}}$,
  receiving $E_{\text{EUR}/\text{USD}}$ USD.
  \item Sell USD for JPY at rate $E_{\text{USD}/\text{JPY}}$,
  receiving $E_{\text{EUR}/\text{USD}} \cdot
  E_{\text{USD}/\text{JPY}}$ JPY.
  \item Sell JPY for EUR at rate $E_{\text{JPY}/\text{EUR}}$,
  receiving the round-trip total
  $E_{\text{EUR}/\text{USD}} \cdot E_{\text{USD}/\text{JPY}}
  \cdot E_{\text{JPY}/\text{EUR}}$ EUR.
\end{enumerate}
The three trades must execute as simultaneously as the matching
engines permit, because rates move while legs are pending. Real
desks send all three orders in a single market message or use
guaranteed-fill order types.

\begin{figure}[ht]
\centering
\begin{tikzpicture}[
  cur/.style={circle, draw, minimum size=1.6cm, text centered,
    font=\small\bfseries, fill=yellow!20},
  arr/.style={-Stealth, thick, blue!65},
  seg/.style={thick, blue!65},
]
\node[cur] (eur) at (90:4.0cm)  {EUR};
\node[cur] (usd) at (210:4.0cm) {USD};
\node[cur] (jpy) at (330:4.0cm) {JPY};
% Diagonals are drawn as two collinear halves meeting at the chord midpoints
% (150:2.0cm) and (30:2.0cm) so neither edge sweeps across the centre box.
\draw[seg] (eur) -- node[pos=1, left=5pt, font=\small]{$E_{\mathrm{EUR/USD}}$} (150:2.0cm);
\draw[arr] (150:2.0cm) -- (usd);
\draw[arr] (usd) -- node[below=5pt, font=\small]{$E_{\mathrm{USD/JPY}}$} (jpy);
\draw[seg] (jpy) -- node[pos=1, right=5pt, font=\small]{$E_{\mathrm{JPY/EUR}}$} (30:2.0cm);
\draw[arr] (30:2.0cm) -- (eur);
\node[draw, rounded corners=3pt, fill=red!8, font=\footnotesize, align=center,
  inner sep=4pt] at (0,0)
  {Arbitrage when\\[2pt]
   $E_{\mathrm{EUR/USD}}$\\[-2pt]$\times\;E_{\mathrm{USD/JPY}}$\\[-2pt]$\times\;E_{\mathrm{JPY/EUR}}\neq 1$};
\end{tikzpicture}
\caption{Triangular FX arbitrage. Starting with one EUR, convert to USD at
$E_{\mathrm{EUR/USD}}$, then to JPY at $E_{\mathrm{USD/JPY}}$, then back to
EUR at $E_{\mathrm{JPY/EUR}}$. If the product of the three rates exceeds one,
a riskless profit exists. Competitive markets enforce the no-arbitrage
constraint that the product equals one. In practice the window is
milliseconds wide and requires co-located systems to exploit.}
\label{fig:triangular_arb}
\end{figure}

\textbf{P\&L.} Per unit of starting EUR, $\text{profit} =
\Delta$. Locked the moment all three legs fill; no holding
period, no subsequent price exposure.

\paragraph{Worked numerical example.} Suppose at some instant
the three rates are
\begin{align*}
  E_{\text{EUR}/\text{USD}} &= 1.08, \\
  E_{\text{USD}/\text{JPY}} &= 148, \\
  E_{\text{JPY}/\text{EUR}} &= 0.0063.
\end{align*}
Then the round-trip is
\begin{equation*}
  \underbrace{1.08}_{\text{EUR}\to\text{USD}}
  \;\times\;
  \underbrace{148}_{\text{USD}\to\text{JPY}}
  \;\times\;
  \underbrace{0.0063}_{\text{JPY}\to\text{EUR}}
  \;=\; 1.006992,
\end{equation*}
so $\Delta \approx 0.006992$, about $0.70\%$ per round trip:
seven thousand EUR on one million, before fees and the cost of
executing fast enough that no leg moves between trades.

The JPY/EUR quote of $0.0063$ implies an EUR/JPY of $\approx
158.73$. The no-arb inverse of the first two rates is
$1/(1.08 \cdot 148) \approx 0.006257$; the quoted $0.0063$
exceeds that by exactly what the round trip captures. In
general: fix any two rates and solve for the third; any
deviation is the locked-in profit.

\textbf{Why it usually doesn't exist.} The major-pair FX market
is the most heavily arbitraged in the world. EBS, Reuters, and
bank-internal ECNs run continuous triangular consistency
monitors, and any basis-point opening closes in microseconds.
\citeBouchaud{} report that on liquid majors the indicator lives
almost entirely within the constituent bid--ask spreads:
apparent gaps are not capturable after crossing the spread on
each leg.

\textbf{When it does.} Three families. (i) \emph{Fragmented
liquidity}: a small ECN quoting all three pairs with fewer
participants drifts off consensus, and a desk plugged into both
the ECN and consensus venues captures the gap. (ii)
\emph{News-driven dislocations}: news hitting EUR/USD a
millisecond before EUR/JPY leaves the triangle inconsistent for
the propagation delay; the FX latency arms race is a race to
close that gap. (iii) \emph{Exotic pairs}: EUR $\to$ TRY $\to$
RUB $\to$ EUR is much less competed; the indicator opens further
and stays open longer in thinner pairs.

\begin{prosperitybox}
Round~$1$'s manual challenge in Prosperity~$3$ was a literal
triangular FX exercise: a $4 \times 4$ conversion matrix between
\texttt{Snowballs}, \texttt{Pizza}, \texttt{Silicon Nuggets},
\texttt{SeaShells}, with the task of finding a sequence of at
most five trades maximising final \texttt{SeaShells}. The
Frankfurt Hedgehogs enumerated all sequences and found the
optimum \texttt{SeaShells} $\to$ \texttt{Snowballs} $\to$
\texttt{Silicon Nuggets} $\to$ \texttt{Pizza} $\to$
\texttt{Snowballs} $\to$ \texttt{SeaShells}, yielding $\approx
8.9\%$ per round trip \citep{frankfurtHedgehogs2025imcP3}. This
is the brute-force version of triangular arbitrage: enumerate
cycles in the currency graph, take the one whose product of edge
weights most exceeds one. In the wild the market enumerates
continuously and closes cycles within microseconds; in
Prosperity the matrix was static for the round, making the
``arbitrage'' a one-shot optimisation. The pedagogical point
survives: this is triangular FX arbitrage in miniature.
\end{prosperitybox}

\section{Locational arbitrage}
\label{sec:ch13_locational}

Same trade, but across a pair of \emph{venues} for the same
asset rather than around a currency triangle.

\textbf{Observable.} Identical contract, identical settlement,
identical quality, different prices on two exchanges. Classical
example: BTC at \$$50{,}000$ on exchange~A and \$$50{,}050$ on
exchange~B. Equity examples include dual-listed stocks on
NYSE/Toronto or E-mini contracts on CME vs alternative venues.
The observable is $\text{ask}_{A} < \text{bid}_{B}$ (crossable
in the lift-A, hit-B direction).

\textbf{Trade.} Buy on the cheaper venue, sell on the dearer, in
equal size. If the asset is fungible and venues share custody
(or the trader holds inventory at both), the position closes
immediately at the gross spread. If the asset must be
\emph{moved} between venues (as for crypto, where transfer
takes minutes to hours and incurs network fees), it is not
strictly an arbitrage during transit: price can move and the
trader wears the risk.

\textbf{P\&L.} Per unit, gross $= \text{bid}_{B} -
\text{ask}_{A}$. Net subtracts (i) exchange fees on both venues,
(ii) transfer cost, (iii) borrow cost if shorting one side,
(iv) slippage beyond top of book. In the BTC example gross is
\$$50$ per coin; net is typically a small fraction of gross.

\textbf{Why it usually doesn't exist.} Automated cross-venue
bots continuously monitor every pair of liquid venues and close
any gap beyond net costs in subsecond time. On BTC/USD between
Coinbase and Binance the gap exceeds joint fees for less than a
fraction of a percent of trading time, and closes within
seconds when it does. The closing \emph{is} the arb; its
existence maintains convergence.

\textbf{When it does.} Three families:
\begin{itemize}
  \item \textbf{Regulatory fragmentation.} Pre-$2018$ Japan vs
  US crypto: capital controls made moving BTC expensive or
  illegal, and BTC traded at a $5$--$10\%$ premium on Japanese
  venues for sustained periods. The Coinbase / bitFlyer spread
  was a multi-month feature, because the closing trade was
  blocked by regulation.
  \item \textbf{Latency asymmetry.} CME E-mini S\&P and
  less-liquid alternative venues briefly dislocate during fast
  moves: a market-maker quoting both from one feed updates one
  before the other. Co-located firms capture the gap in a few
  hundred microseconds.
  \item \textbf{Transient liquidity imbalance.} A large
  aggressive order drives the local bid--ask off consensus; the
  asymmetry persists for seconds before cross-venue MMs
  re-equalise. A bot positioned on both venues lifts the cheap
  side and hits the dear side during the window.
\end{itemize}

\begin{prosperitybox}
Round~$4$ of Prosperity~$3$ introduced
\texttt{MAGNIFICENT\_MACARONS}, tradeable on the local order
book and via a fixed-fee external ``conversion'' channel
exposing an external bid/ask adjusted for import/export and
transport fees \citep{frankfurtHedgehogs2025imcP3}. Whenever
local bid exceeded external ask (after fees), or local ask fell
below external bid, a profitable conversion was available:
textbook locational arbitrage with the local book as one venue
and the external channel as the other. The Frankfurt Hedgehogs
report that over $10{,}000$ timesteps with a $10$-unit per-tick
conversion cap, the trade could theoretically yield up to
$300{,}000$ SeaShells per round, with realistic captures of
$130{,}000$--$160{,}000$. Their implementation placed limit
sells on the local book at $\lfloor\text{externalBid} +
0.5\rfloor$ each tick (the maximum price that still attracted
local taker bots) and converted resulting inventory through
the external channel. Locational arbitrage can be a stable,
high-throughput source of profit when the closing mechanism is
rate-limited.
\end{prosperitybox}

\section{Cash-and-carry arbitrage}
\label{sec:ch13_cash_carry}

Cash-and-carry links spot and futures on the same asset. Its
no-arbitrage relation $F = S e^{(r-y)(T-t)}$ is one of the
half-dozen most fundamental formulas in derivative pricing.

\subsection{The cost-of-carry relation, derived from no-arbitrage}
\label{sec:ch13_cost_of_carry_derivation}

Consider a non-dividend-paying asset with spot $\Spx$, risk-free
rate $\rrf$ (borrow = lend), and a futures contract with
delivery $T$ and price $F$. We derive $F$ by exhibiting a
replicating portfolio.

\emph{Portfolio~A}: at $t=0$, borrow $\Spx$ in cash, buy one
unit at spot, and enter a short futures with delivery price $F$.
Setup is cashless (borrowed cash pays for the asset; the short
futures requires no upfront premium beyond margin, treated as
zero-cost). At time $T$:
\begin{itemize}
  \item Deliver into the short futures for $F$ cash.
  \item Repay the loan, grown to $\Spx e^{\rrf T}$ under
  continuous compounding.
\end{itemize}
The terminal cashflow $F - \Spx e^{\rrf T}$ is deterministic
(both legs locked at $t=0$), so no-arbitrage forces it to zero:
\begin{equation}
  F \;=\; \Spx \, e^{\rrf T},
  \label{eq:cost_of_carry_no_div}
\end{equation}
which is the no-dividend, no-storage-cost version of the
cost-of-carry formula.

If $F > \Spx e^{\rrf T}$, Portfolio~A delivers positive cashflow
at zero cost: borrow cash, buy spot, short futures, deliver,
pocket the difference (\emph{cash-and-carry arbitrage}). If
$F < \Spx e^{\rrf T}$, the reverse trade (short spot, lend
cash, long futures, take delivery) is the arbitrage
(\emph{reverse cash-and-carry}).

\paragraph{Adding a yield.} If the asset pays a continuous yield
at rate $y$ (dividend yield for equities, coupons for bonds,
convenience yield for commodities useful to hold), Portfolio~A's
holding leg receives extra $\Spx(e^{yT}-1)$. Re-running the
argument:
\begin{equation}
  F \;=\; \Spx \, e^{(\rrf - y) T},
  \label{eq:cost_of_carry}
\end{equation}
which is the general cost-of-carry forward formula.

\begin{keyfactbox}[Cost-of-carry forward formula]
\label{kf:cost_of_carry}
For an asset with spot price $\Spx$, risk-free rate $\rrf$, and
continuous yield $y$ (dividend yield for equities, convenience
yield net of storage for commodities), the no-arbitrage
forward / futures price for delivery at time $T$ is
\begin{equation}
  F \;=\; \Spx \, e^{(\rrf - y) T}.
\end{equation}
A \emph{cash-and-carry arbitrage} exists when the market
forward $F_{\text{market}}$ differs from this fair value:
\begin{itemize}
  \item $F_{\text{market}} > \Spx e^{(\rrf-y)T}$: long spot,
  short futures, finance at $\rrf$, deliver at $T$. Locks in
  $F_{\text{market}} - \Spx e^{(\rrf-y)T}$ per unit.
  \item $F_{\text{market}} < \Spx e^{(\rrf-y)T}$: short spot,
  long futures, lend at $\rrf$, take delivery at $T$. Locks in
  $\Spx e^{(\rrf-y)T} - F_{\text{market}}$ per unit.
\end{itemize}
\end{keyfactbox}

\subsection{Worked numerical example}
\label{sec:ch13_carry_example}

Consider an equity index with spot $\Spx = 100$, risk-free rate
$\rrf = 5\%$, dividend yield $y = 1\%$, and a futures contract
expiring in $T = 0.5$ years. The fair forward is
\begin{equation*}
  F \;=\; 100 \cdot e^{(0.05 - 0.01) \cdot 0.5}
  \;=\; 100 \cdot e^{0.02}
  \;=\; 100 \cdot 1.02020134\ldots
  \;\approx\; 102.02.
\end{equation*}
Suppose the market quotes the futures at $F_{\text{market}} =
103$. Then the cash-and-carry arbitrage is open at
\begin{equation*}
  103 - 102.02 \;\approx\; 0.98 \text{ per unit}.
\end{equation*}
At $t=0$: borrow $100$ at $5\%$, buy one index unit, sell one
futures at $103$. At $T=0.5$: receive dividends
$100(e^{0.005}-1) \approx 0.5013$ (continuous compounding),
deliver for $103$, repay loan $100 \cdot e^{0.025} \approx
102.5316$. Net at $T$:
\begin{equation*}
  \underbrace{103}_{\text{delivery}}
  \;+\;
  \underbrace{0.5013}_{\text{dividend}}
  \;-\;
  \underbrace{102.5316}_{\text{loan repayment}}
  \;=\;
  0.9697 \text{ per unit},
\end{equation*}
matching the $\sim 0.98$ headline up to compounding rounding.
The arb locks the moment loan, spot, and futures legs are in
place.

\textbf{Why it usually doesn't exist.} Index-futures-vs-spot
bases are watched continuously by every relative-value desk.
$\rrf$ is well-known, $y$ is forecast from public dividend
schedules, and $F_{\text{market}} - \Spx e^{(\rrf - y)T}$
typically lives below joint round-trip cost. The basis is the
cleanest and most heavily competed trade in equities.

\textbf{When it does.} Cash-and-carry opens when one formula
input moves faster than the market can price it in:
\begin{itemize}
  \item \textbf{Dividend-date dislocations.} On a large ex-div
  morning the cash market drops by the dividend; if futures
  lags, the basis opens by a comparable amount.
  \item \textbf{Forced unwinds.} A margin-called leveraged fund
  dumps spot; spot drops below fair value while futures stays
  anchored, opening a reverse cash-and-carry.
  \item \textbf{Commodity storage constraints.} For storage-
  constrained commodities (April~$2020$ WTI), the convenience
  yield $y$ turns large and negative, and the futures curve
  contangoes beyond rates-implied fair value. The arbitrageur
  cannot find a tank to store physical oil: the ``free''
  part of the free lunch is not free.
  \item \textbf{Funding-rate spikes.} If $\rrf$ jumps
  (year-end repo squeezes), $\Spx e^{\rrf T}$ jumps with it; if
  the futures curve does not follow, the basis opens.
\end{itemize}

Hull \citeHull{} (Ch~$5$) is the standard reference.

\section{Put-call parity arbitrage}
\label{sec:ch13_put_call_parity}

Put-call parity is the no-arbitrage relation between a European
call, a European put, and the underlying. It is the only formula
here that requires knowing what calls and puts are; we state the
payoffs and derive parity directly. Full treatment:
\cref{ch:options_payoffs}.

\subsection{Setup}
\label{sec:ch13_parity_setup}

A \emph{European call} with strike $\Kstrike$ and maturity
$\Tmat$ on underlying $\Spx$ pays $\max(\Spx_T - \Kstrike, 0)$
at $\Tmat$; its price today is $\Call$. A \emph{European put}
with the same strike and maturity pays $\max(\Kstrike - \Spx_T,
0)$ and has price $\Put$. The risk-free rate is $\rrf$. We
assume no dividends between today and $\Tmat$; the dividend
correction is in the pitfall box below.

\subsection{Derivation from no-arbitrage}
\label{sec:ch13_parity_derivation}

Consider the portfolio: long one call, short one put, lend
$\Kstrike e^{-\rrf \Tmat}$ (equivalently, a zero-coupon bond
paying $\Kstrike$ at $\Tmat$). Cost today:
\begin{equation*}
  V_0 \;=\; \Call - \Put + \Kstrike e^{-\rrf \Tmat}.
\end{equation*}
At $\Tmat$ the portfolio value depends on $\Spx_T$:

\textbf{Case~1: $\Spx_T > \Kstrike$.} Call pays $\Spx_T -
\Kstrike$; short put expires worthless; bond pays $\Kstrike$.
Total:
\begin{equation*}
  V_T \;=\; (\Spx_T - \Kstrike) + 0 + \Kstrike \;=\; \Spx_T.
\end{equation*}

\textbf{Case~2: $\Spx_T \leq \Kstrike$.} Call expires worthless;
short put owes $\Kstrike - \Spx_T$; bond pays $\Kstrike$. Total:
\begin{equation*}
  V_T \;=\; 0 - (\Kstrike - \Spx_T) + \Kstrike \;=\; \Spx_T.
\end{equation*}

Both cases give $V_T = \Spx_T$, so the portfolio replicates one
share and its cost today must equal the spot price:
\begin{equation*}
  \Call - \Put + \Kstrike e^{-\rrf \Tmat} \;=\; \Spx.
\end{equation*}
Rearranging gives the famous put-call parity identity.

\begin{keyfactbox}[Put-call parity]
\label{kf:put_call_parity}
For European call $\Call$ and put $\Put$ with the same strike
$\Kstrike$ and maturity $\Tmat$, on a non-dividend-paying
underlying $\Spx$ at risk-free rate $\rrf$,
\begin{equation}
  \Call - \Put \;=\; \Spx - \Kstrike e^{-\rrf \Tmat}.
  \label{eq:put_call_parity}
\end{equation}
The identity holds at every instant before maturity by
no-arbitrage alone, with no dynamic model and no distributional
assumption, only the absence of free lunches.
\end{keyfactbox}

\begin{intuitionbox}
Picture the payoffs. Long call: zero below $\Kstrike$, slopes
up at $45^\circ$ above. Short put: zero above $\Kstrike$,
slopes down at $45^\circ$ below. Stacked, they form a single
$45^\circ$ line, a synthetic long forward struck at
$\Kstrike$. A forward to buy at $\Kstrike$ is worth $\Spx -
\Kstrike e^{-\rrf \Tmat}$ by cost-of-carry, hence $\Call - \Put
= \Spx - \Kstrike e^{-\rrf \Tmat}$. Parity and cost-of-carry are
two views of the same no-arbitrage skeleton.
\end{intuitionbox}

\subsection{The arbitrage trade}
\label{sec:ch13_parity_trade}

\textbf{Observable.} A market in which $\Call$, $\Put$, $\Spx$
violate \eqref{eq:put_call_parity}. Two cases:
\begin{itemize}
  \item $\Call - \Put > \Spx - \Kstrike e^{-\rrf \Tmat}$: the
  call--put combination is overvalued relative to the
  synthetic-forward right-hand side.
  \item $\Call - \Put < \Spx - \Kstrike e^{-\rrf \Tmat}$: the
  call--put combination is undervalued.
\end{itemize}

\textbf{Trade.} Overvalued case (the other is symmetric): sell
call, buy put, buy underlying, borrow $\Kstrike e^{-\rrf
\Tmat}$. Cash inflow at $t=0$ is
\begin{equation*}
  +\Call - \Put - \Spx + \Kstrike e^{-\rrf \Tmat}
  \;=\; (\Call - \Put) - (\Spx - \Kstrike e^{-\rrf \Tmat}),
\end{equation*}
the parity violation, positive by hypothesis. At maturity the
four legs cancel by the two-state argument, so the violation is
pocketed and nothing is owed later.

\textbf{P\&L.} The parity violation, locked at $t=0$.

\textbf{Why it usually doesn't exist.} Listed options market
makers enforce parity continuously: their quoting algorithms
compute, per (strike, maturity) pair, the call implied by the
put (or vice versa) via parity, and refuse to let quotes drift
more than a tick from consistency. A visible parity gap on a
liquid name is usually a stale quote.

\textbf{When it does.} Parity opens when the underlying leg
becomes non-trivial:
\begin{itemize}
  \item \textbf{Hard-to-borrow underlyings.} For heavily-shorted
  small-caps or squeezed meme stocks, the short-underlying leg
  carries real borrow cost, modelled as a positive $y$ in the
  cost-of-carry relation. Assuming $y=0$ produces fake
  violations that are not capturable.
  \item \textbf{Discrete dividends.} A dividend between $t$ and
  $\Tmat$ goes to the underlying holder but not to the synthetic
  combination; the formula subtracts $PV$(div). See the pitfall
  box.
  \item \textbf{Corporate actions.} Mergers, spin-offs, special
  dividends, splits disrupt the options chain while the exchange
  recalculates specs; parity can drift during the adjustment.
  \item \textbf{Early-exercise on American options.} Parity
  holds strictly only for European options; American options
  give a pair of inequalities, one of the six no-arb relations
  in \citep{damodaranOptionsArbitrage}.
\end{itemize}

\subsection{Worked numerical example}
\label{sec:ch13_parity_example}

Take $\Spx = 100$, $\Kstrike = 100$, $\rrf = 5\%$, $\Tmat = 1$
year, and a non-dividend-paying underlying. The right-hand side
of parity is
\begin{equation*}
  \Spx - \Kstrike e^{-\rrf \Tmat}
  \;=\; 100 - 100 \cdot e^{-0.05}
  \;=\; 100 \cdot (1 - e^{-0.05})
  \;=\; 100 \cdot 0.04877\ldots
  \;\approx\; 4.877.
\end{equation*}
So put-call parity says $\Call - \Put \approx 4.877$.

Suppose the market quotes $\Call = 8$, $\Put = 2$. Then $\Call -
\Put = 6$ exceeds the parity-implied $4.877$ by $6 - 4.877
\approx 1.123$ per pair, which the arbitrageur locks by selling
the call, buying the put, buying the stock at $100$, and
borrowing $100 e^{-0.05} \approx 95.123$. Net inflow $t=0$ is
$1.123$; at maturity the four legs cancel.

\begin{pitfallbox}[Parity breaks with discrete dividends]
\label{pf:parity_dividends}
$\Call - \Put = \Spx - \Kstrike e^{-\rrf \Tmat}$ assumes no
dividends. Real equities pay discrete dividends to whoever
holds the share \emph{on the ex-dividend date}, not to the
holder of the synthetic call--put combination. Corrected
parity:
\begin{equation*}
  \Call - \Put \;=\; \Spx - PV(\text{div}) - \Kstrike e^{-\rrf \Tmat},
\end{equation*}
where $PV$(div) is the present value of dividends scheduled
between $t$ and $\Tmat$. Forgetting this term is a common
scanner failure: every quarter, dividend-paying stocks report a
dividend-sized ``violation'' that is entirely fake. Always strip
the dividend stream. The same correction applies to
\eqref{eq:cost_of_carry} for equity indices with non-zero
$y$: index versions use a continuous-yield approximation,
single-name equities need the discrete version.
\end{pitfallbox}

The Damodaran catalogue \citep{damodaranOptionsArbitrage}
develops six no-arb relations on options (intrinsic-value
bounds, strike and maturity monotonicity, and parity), each
derived by the same pattern: two-state payoff at expiry, show
portfolio payoff non-negative in both states, conclude cost
today non-positive. Hull \citeHull{} (Ch~$5$) is the standard
textbook treatment.

\section{Index arbitrage}
\label{sec:ch13_index}

Index arbitrage is cost-of-carry on an index futures contract,
where the underlying is a basket (for the S\&P~$500$, five
hundred stocks in declared proportions). The fair forward comes
from \eqref{eq:cost_of_carry} with $\Spx$ the weighted
constituent sum and $y$ the weighted-average dividend yield.

\textbf{Observable.} Index futures away from
\begin{equation*}
  F^*_{\text{index}} \;=\; \Spx_{\text{SPX}} \cdot e^{(\rrf - y)(T - t)},
\end{equation*}
with $\Spx_{\text{SPX}}$ the $500$-stock synthetic at wall-mid
prices, $\rrf$ the overnight rate (fed funds or SOFR), $y$ the
index dividend yield.

\textbf{Trade.} If $F_{\text{market}} > F^*_{\text{index}}$:
short futures, buy all $500$ constituents in index proportions.
If $F_{\text{market}} < F^*_{\text{index}}$: reverse. Hold to
expiry; deliver (or cover). Dedicated ``program trading'' or
``index arb'' desks exist to run this, the largest single
trade by volume in equities.

\textbf{P\&L.} Per index unit, $F_{\text{market}} -
F^*_{\text{index}}$, times contract notional and size.

\textbf{Why it usually doesn't exist.} Index arb is the most
heavily competed trade in the world. The basis typically lives
inside the joint bid--ask of the futures and the cheapest
tracking portfolio. Convergence is so tight that the basis is
sometimes read as a real-time benchmark for $\rrf$: given
observable $F_{\text{market}}$ and $\Spx_{\text{SPX}}$,
\eqref{eq:cost_of_carry} is inverted for the market's funding
cost.

\textbf{When it does.} Index arb opens at moments of
constituent-level chaos:
\begin{itemize}
  \item \textbf{Index-rebalance events.} When S\&P changes
  membership on a known date, index funds must track the new
  composition; the basis can dislocate during the recompute
  window.
  \item \textbf{Quadruple witching.} Once a quarter, futures,
  options, and underlyings expire together, triggering massive
  program-trading flow and occasional dislocations beyond joint
  transaction cost.
  \item \textbf{Fast-market stress.} During extreme vol
  (March~$2020$ COVID, August~$2024$ Yen unwind), the $500$-stock
  basket does not track futures in lock-step; the basis can blow
  out by percentage points for tens of minutes.
\end{itemize}

\paragraph{Connection to basket arbitrage.} Mathematically
identical to \cref{ch:basket_etf_arb}: both compare a synthetic
weighted-constituent value to a single wrapper quote (futures or
ETF). The structural difference is the closing mechanism:
creation/redemption for ETFs, physical delivery at maturity for
futures.

\section{Calendar arbitrage}
\label{sec:ch13_calendar}

Calendar arbitrage exploits inconsistent pricing between two
contracts on the same underlying with different expiries. The
common venue is options, where the implied-vol term structure
occasionally inverts by more than the calendar no-arb bound
permits.

\textbf{Observable.} Two options on the same underlying, same
strike, maturities $T_1 < T_2$. The no-arb bound says the
longer-dated option is at least as valuable (it retains option
value after the shorter expires): for European calls,
$\Call(T_2) \geq \Call(T_1)$. A violation is the observable, or
equivalently a term-structure inversion where short-dated vol
exceeds long-dated vol by more than the bound permits. Hull
\citeHull{} (Ch~$17$) covers calendar spreads.

\textbf{Trade.} Buy the cheap (long-dated), sell the expensive
(short-dated), hold to short-dated expiry. The short expires;
the trader is left holding the long, worth at least what the
short was, plus the captured violation.

\textbf{P\&L.} The violation net of joint transaction cost. This
is a true Type~I arbitrage only if the long is held to its own
expiry and residual value works out; in practice desks unwind
earlier and treat it as a relative-value position.

\textbf{Why it usually doesn't exist.} Options market makers
smooth the vol surface by construction: quoting algorithms
enforce monotonicity and convexity in strike and maturity, and
errors correct within seconds. The calendar bound is one of the
six monotonicity conditions in
\citep{damodaranOptionsArbitrage}.

\textbf{When it does.} Around scheduled events between the two
maturities (earnings, central bank meetings, FDA decisions)
where short-dated implied vol spikes for event risk while
long-dated does not move much (the event is a small fraction of
its remaining variance). Biotech FDA decisions are the classic
case: short-dated implied vol can grow by an order of magnitude
in the days before announcement.

A second source is liquidity: for thin long-dated markets
(exotic commodity options, deep-OTM LEAPS), the long quote goes
stale relative to the short, inverting the term structure as a
quoting artefact. The detection logic is the same though the
capture looks more like market making.

\section{ETF creation/redemption arbitrage}
\label{sec:ch13_etf_cr}

The mechanism is in \cref{ch:basket_etf_arb}; this section
places the trade in the zoo and draws the line between it and
the basket-spread stat-arb of the same chapter.

\textbf{Observable.} ETF price diverges from synthetic basket
value by more than the creation/redemption fee plus joint
transaction cost. With ETF mid $P_{\text{ETF}}$ and synthetic
NAV $P_{\text{NAV}}$:
\begin{equation*}
  |P_{\text{ETF}} - P_{\text{NAV}}|
  \;>\;
  \text{creation fee} + \text{transaction costs}.
\end{equation*}

\textbf{Trade.} If $P_{\text{ETF}} > P_{\text{NAV}}$: the AP
shorts the ETF, buys the basket, uses the basket to
\emph{create} new ETF shares with the issuer, delivers against
the short. If $P_{\text{ETF}} < P_{\text{NAV}}$: the AP buys
the ETF, redeems with the issuer for the basket, sells the
basket.

\textbf{P\&L.} $|P_{\text{ETF}} - P_{\text{NAV}}|$ minus
creation fee minus joint transaction cost, locked per creation
unit ($\sim 50{,}000$ shares).

\textbf{Why it usually doesn't exist.} APs run automated
basis-monitoring on every liquid ETF and are contractually
incentivised to keep the basis tight. On SPY, IVV, VOO the
basis rarely exceeds a fraction of a basis point intraday, and
openings close in seconds.

\textbf{When it does.} Three families:
\begin{itemize}
  \item \textbf{Creation fee changes.} A fee increase widens the
  no-arb band; a trader who knows the new schedule before the
  basis adjusts captures the transition.
  \item \textbf{Settlement lag.} For ETFs with international
  constituents (EM ETFs spanning $20$ time zones), constituent
  prices are stale when home markets are closed; the basis is a
  noisy stale-vs-current mix. Sophisticated APs use fair-value
  models; traders who model time zones better than consensus
  find real openings.
  \item \textbf{AP outages.} When an AP is unavailable (trading
  system down, credit line frozen), the closing mechanism is
  gone. The basis can blow out by percent until another AP picks
  up. The ``flash crash'' episodes are the dramatic cases: ETF
  prices collapse relative to NAV when AP-side and exchange-side
  liquidity vanish together.
\end{itemize}

The Prosperity~$3$ \texttt{PICNIC\_BASKET} trade in
\cref{ch:basket_etf_arb} is \emph{not} an instance: the
simulator has no creation/redemption mechanism; basket-to-NAV
convergence is a property of the data-generating process, not a
closed loop, so the trade has no zero-risk completion. The next
section is about that distinction.

\section{Statistical arbitrage: not a true arbitrage}
\label{sec:ch13_stat_arb}

The seven preceding sections cover \emph{true arbitrage}:
deterministic, non-negative-in-every-state by construction, with
a no-arb argument that pins fair value down.

The strategies of \cref{ch:mean_reversion_ou},
\cref{ch:pairs_cointegration}, and \cref{ch:basket_etf_arb}
(outside the literal ETF c/r mechanism above) are different in
kind. They share the word ``arbitrage'' for historical/marketing
reasons but are not arbitrages in the sense of
\cref{def:arbitrage}. A statistical arbitrage:
\begin{itemize}
  \item Has positive \emph{expected} payoff under an assumed
  data-generating process.
  \item Has non-zero variance of realised payoff.
  \item Can lose in any individual trade and over any finite
  window.
  \item Requires an assumption about the persistence of a
  statistical relationship, with no contractual no-arb
  guarantee enforcing it.
\end{itemize}

The OU $z$-score rule of \cref{ch:mean_reversion_ou} is
positive-EV under the assumption that the spread follows OU
with the calibrated parameters; if parameters change, the rule
can be wrong for arbitrarily long and the position can lose
arbitrarily much before convergence. The cointegration tests of
\cref{ch:pairs_cointegration} are statistical inferences that
can be correct in expectation and wrong on the realised path.
The \texttt{PICNIC\_BASKET1} spread trade works in Prosperity
because the quoting bots are programmed to mean-revert toward
synthetic value (a property of the simulation, not the
universe); a different parameterisation breaks it.

Sizing follows from the distinction. A true arb is sized to
maximum capacity, because there is no downside; the only
constraint is broker-deployable principal. A stat arb must be
sized via \cref{ch:kelly_risk}: by the ratio of expected return
to variance, with full Kelly as an upper bound and some
fraction thereof in practice. Sizing a stat arb as a true arb
is unbounded risk in disguise.

\begin{pitfallbox}[Statistical arbitrage is not arbitrage]
\label{pf:stat_arb_is_not}
``Arbitrage'' in ``statistical arbitrage'' is a marketing
label from the late $1980$s/early $1990$s, not the technical
term of \cref{def:arbitrage}. A stat arb has positive expected
return, not deterministic, and it can and does lose. LTCM
($1998$, \citep{lowenstein2000ltcm}), the $2007$ quant quake
that vaporised equity-market-neutral funds in three days, and a
long list of smaller blow-ups reduce to one error: a stat-arb
book sized as if it could not lose, then realising
cross-position-correlated losses far beyond the in-sample
variance estimate.

Two rules:
\begin{enumerate}
  \item Never call a bet an ``arbitrage'' or an arbitrage a
  ``trade''. Vocabulary discipline is the first defence against
  the sizing error.
  \item Size stat arbs by Kelly (\cref{ch:kelly_risk}), never by
  ``fully deploy capital''. The bound is the risk model's worst
  plausible drawdown, not the capital balance.
\end{enumerate}
Apply the five-step template of \cref{kf:arb_template} and ask
at step~(3): \emph{is the P\&L deterministic?} For pairs
trading, non-creatable basket spreads, and OU mean reversion
the answer is no, and the trade moves out of the zoo and into
the strategy bestiary.
\end{pitfallbox}

\begin{gsbox}
Every large bank runs a desk, called \emph{relative value},
\emph{program trading}, or simply \emph{the arb desk}, mandated
to monitor and trade the seven arbitrages here. On a Goldman Sachs
floor in $2025$, a junior strat's morning begins with a
\emph{zoo sweep}: a scripted walk-through of every arbitrage
against intraday prices, with any opening flagged in a Slack
channel the senior traders read over first coffee. On most
mornings every entry returns ``closed by less than round-trip
cost'' and nothing happens. On the rare morning one entry
opens (a parity violation on a hard-to-borrow name, a basis
dislocation around quadruple witching, an index rebalance, an
FDA decision), the trader descends for an hour and the desk
makes a meaningful share of its monthly P\&L. The
patience-and-automation culture flows from this rhythm: daily
expected return is flat, but conditional expected return given
an opening is huge, so the trader set up to capture openings
the moment they appear is the one who eats. The strats
building zoo-sweep monitors are accordingly some of the
quietly important people on the floor: their dashboards are
the institutional memory of when, how, and why each arbitrage
has historically opened.
\end{gsbox}

\begin{historybox}
Systematic study of no-arb relations dates to the early
$1970$s. Black--Scholes ($1973$ JPE) with Merton's companion
\citep{merton1973rational} derived \eqref{eq:put_call_parity}
as a corollary of the replicating-portfolio argument behind the
Black--Scholes equation. \citet{stoll1969relationship} is the
canonical pre-1973 citation for the US-equity parity
inequality, but the modern hierarchical view, with parity as the
simplest of a ladder topped by the Black--Scholes PDE, comes
from the $1973$ papers.

Institutional arbitrage on a trading floor is older. Gustave
Levy ran Goldman Sachs' arb desk from the late $1940$s and was
Senior Partner by $1969$, making his name on \emph{merger
arbitrage}: buying a target at a discount to the deal price,
hedged by shorting the acquirer. Merger arb is not in the zoo
(deals break, and the spread compensates deal-break risk),
but Goldman's modern arb desk traces to Levy's $1950$s practice.

The cautionary half is LTCM, founded $1994$ and bankrupt
$1998$ \citep{lowenstein2000ltcm}. Founded by Meriwether with
Merton and Scholes themselves, LTCM ran statistical arbitrages
dressed as textbook arbs (on-the-run/off-the-run Treasury
spreads, swap spreads, sovereign convergence), all
positive-EV with small normal-regime variance, all
simultaneously losing during the August--September~$1998$
Russian-default flight to quality. A Fed-organised \$$3.6$B
private-sector recapitalisation averted broader collapse. The
lesson: a strategy whose name contains ``arbitrage'' is not, by
that fact alone, an arbitrage.
\end{historybox}

\section{The arb-diagnostic flowchart}
\label{sec:ch13_flowchart}

The chapter ends where it began: with a diagnostic. Given an
observable condition, which arbitrage fits? The table lists the
seven zoo entries with their triggers, trade structures, and
section references. Keep it in arm's reach when encountering an
unfamiliar product and walk down row by row.

\renewcommand{\arraystretch}{1.3}
\begin{center}
\begin{tabular}{p{0.34\linewidth} p{0.34\linewidth} p{0.20\linewidth}}
\toprule
\textbf{If you observe\ldots} & \textbf{\ldots set up} & \textbf{See} \\
\midrule
Three FX rates whose product around a triangle differs from
one. &
Triangular arbitrage: cycle through the three pairs in the
profitable direction. &
\cref{sec:ch13_triangular} \\
\midrule
The same asset quoted at materially different prices on two
venues, after fees. &
Locational arbitrage: buy on the cheap venue, sell on the
dear venue. &
\cref{sec:ch13_locational} \\
\midrule
A futures price away from $\Spx e^{(\rrf - y)(T-t)}$, after
financing and yield. &
Cash-and-carry arbitrage: long spot, short futures (or the
reverse), deliver at maturity. &
\cref{sec:ch13_cash_carry} \\
\midrule
$\Call - \Put$ on a European pair away from $\Spx - \Kstrike
e^{-\rrf \Tmat}$ (after dividends and borrow). &
Put-call parity arbitrage: trade the four legs to lock in the
violation. &
\cref{sec:ch13_put_call_parity} \\
\midrule
Index futures away from the synthetic basket value (after
dividends and rates). &
Index arbitrage: trade the futures against all constituents
in index proportions. &
\cref{sec:ch13_index} \\
\midrule
A short-dated and long-dated option on the same underlying
with relative pricing violating monotonicity. &
Calendar arbitrage: buy the cheap maturity, sell the
expensive, hold to short-leg expiry. &
\cref{sec:ch13_calendar} \\
\midrule
ETF wrapper trading away from synthetic NAV by more than the
creation fee. &
ETF creation / redemption arbitrage: create or redeem with
the issuer to capture the basis. &
\cref{sec:ch13_etf_cr} \\
\midrule
A spread between historically-correlated assets that has
drifted from its long-run mean, with no contractual mechanism
to force convergence. &
\emph{Statistical arbitrage} (NOT a true arb): a positive-EV
bet to be sized via Kelly. &
\cref{sec:ch13_stat_arb}, \cref{ch:kelly_risk} \\
\bottomrule
\end{tabular}
\end{center}

One row per zoo entry, plus the closing row explicitly
excluding stat arb. A reader with this table internalised can
walk into any market, look at available instruments, and
within minutes have an opinion on which arbitrages are worth
monitoring and which are inapplicable. That is the
chapter's operational deliverable.

\section{Key facts to memorise}
\label{sec:ch13_keyfacts}

\begin{itemize}
  \item \textbf{Definition.} An arbitrage is a strategy with
  zero or negative initial cost, non-negative payoff in every
  state of the world, and strictly positive payoff in some
  state. Equivalently, a free lunch.
  \item \textbf{The five-step template.} Observable, trade,
  P\&L, why it usually does not exist, when it does. Every
  entry in the zoo follows this template.
  \item \textbf{Triangular FX.} Three currency pairs whose
  product around a triangle differs from one; trade the cycle.
  Profit per unit is the gap; example $1.08 \cdot 148 \cdot
  0.0063 - 1 \approx 0.70\%$.
  \item \textbf{Locational arbitrage.} Same asset, two venues,
  different prices after fees; buy cheap, sell dear. Opens on
  regulatory fragmentation, latency asymmetry, transient
  imbalance.
  \item \textbf{Cost of carry.} The no-arbitrage forward price
  is $F = \Spx e^{(\rrf - y) T}$. The cash-and-carry arbitrage
  closes any deviation from this; the example with $\Spx =
  100$, $\rrf = 5\%$, $y = 1\%$, $T = 0.5$ gives fair forward
  $\approx 102.02$.
  \item \textbf{Put-call parity.} For European options on a
  non-dividend underlying, $\Call - \Put = \Spx - \Kstrike
  e^{-\rrf \Tmat}$. With $\Spx = \Kstrike = 100$, $\rrf = 5\%$,
  $\Tmat = 1$, parity gives $\Call - \Put \approx 4.877$. Add
  $-PV(\text{div})$ if the underlying pays dividends.
  \item \textbf{Index arbitrage.} Cost-of-carry applied to a
  basket: index futures vs the synthetic value of the
  $500$-stock basket. Most heavily competed trade in equities;
  opens on rebalances, witching days, fast markets.
  \item \textbf{Calendar arbitrage.} Two maturities of the
  same option, monotonicity violated; buy the cheap one, sell
  the expensive one. Opens around scheduled events between
  the two maturities.
  \item \textbf{ETF creation / redemption arbitrage.} ETF
  wrapper price away from synthetic NAV by more than the
  creation fee; create or redeem to close. Distinct from the
  Prosperity \texttt{PICNIC\_BASKET} stat arb of
  \cref{ch:basket_etf_arb}, which has no creation mechanism.
  \item \textbf{Statistical arbitrage is not arbitrage.}
  Positive expected value, non-zero variance, can and does
  lose. Size via Kelly (\cref{ch:kelly_risk}), not by capital
  deployment.
  \item \textbf{The desk rhythm.} Expected open arbitrages per
  day is small; the relative-value desk detects openings
  automatically the moment they appear. The
  patience-and-automation culture follows from the
  conditional-vs-unconditional expected-return asymmetry.
\end{itemize}
